\begin{filecontents*}[overwrite]{verify_hsz_third_coefficient.py}
"""Exact symbolic checks accompanying the HSZ third-coefficient proof.
Requires Python 3 and SymPy. No numerical search is used as a proof.
Run: python verify_hsz_third_coefficient.py
"""
import json
import sympy as sp
results = []
ns = {}
exec('"""Exact identity checks for the one-boundary-root entropy proof.\nNo numerical search or numerical inequality certification is performed.\nAll strict signs and covariance arguments are analytic in the audit note.\n"""\nfrom pathlib import Path\nimport hashlib\nimport json\nimport sympy as sp\nq,s,v,u,x,z=sp.symbols(\'q s v u x z\', real=True)\nd=2*q-1\nE=(s*s-2*v)*(d*v*v+1)-2*s*(d*v+1)*v-2*q*v*v\nchecks=[]\ndef same(name,left,right):\n    residual=sp.factor(sp.together(left-right))\n    assert residual==0,(name,residual)\n    checks.append(name)\nE0=sp.expand(E.subs(q,0));E1=sp.expand(E.subs(q,1))\nsame(\'affineness in q\',E,(1-q)*E0+q*E1)\na=1-u\nsame(\'negative v endpoint q0\',E0.subs(v,-u),(1+u)*(a*s*s+2*u*(s+a)))\nsame(\'negative v endpoint q1\',E1.subs(v,-u),(1+u*u)*s*s+2*u*a*s+2*u*(1-u+u*u))\nsame(\'negative v q1 positive remainder\',E1.subs(v,-u)-2*u*u,(1+u*u)*s*s+2*u*a*(s+a))\nsame(\'mass after s=-x\',q-sp.Rational(1,2)*x*(2*q-1+1/v),q*(1-x)-x*(1-v)/(2*v))\nX=2*v/(1+v)\nsame(\'nonpositive s endpoint q0\',E0.subs(s,-x),(1-v)*((1+v)*x*x+2*v*x-2*v*(1+v)))\nsame(\'nonpositive s endpoint q1\',E1.subs(s,-x),(1+v*v)*x*x+2*v*(1+v)*x-2*v*(1+v+v*v))\nsame(\'q0 endpoint at X\',E0.subs(s,-X),-2*v*(1-v)**3/(1+v))\nsame(\'q1 endpoint at X\',E1.subs(s,-X),-2*v*(1-v)**2*(v*v+v+1)/(1+v)**2)\nA=(1+v)*s*s-2*v*s-2*v*(1+v)\nB=s*s-2*s-2*v-1\nsame(\'A B decomposition\',E,(1-v)*A+2*q*v*v*B)\nT=4*v/(1+v)\nsame(\'A at threshold\',A.subs(s,T),-2*v*(1-v)**2/(1+v))\nP=(1-v)**2+s*s-2*s*(1+v)*x+4*v*x*x\nsame(\'quadratic boundary modulus\',P.subs(x,(z+1/z)/2),(1-s*z+v*z*z)*(1-s/z+v/z**2))\nsame(\'derivative positivity certificate\',-sp.diff(P,x),2*(1+v)*(s-T)+8*v*(1-x))\n# Beta integral of (1-cos theta)^a under normalized circle measure.\nbeta_exponent=sp.symbols(\'beta_exponent\',positive=True)\nJ=2**beta_exponent*sp.gamma(beta_exponent+sp.Rational(1,2))/(sp.sqrt(sp.pi)*sp.gamma(beta_exponent+1))\nsame(\'base probability mass\',J.subs(beta_exponent,1),1)\nsame(\'base entropy\',sp.expand_func(sp.diff(J,beta_exponent).subs(beta_exponent,1)),1-sp.log(2))\n', ns)
results.extend('one-boundary' + ': ' + name for name in ns['checks'])
ns = {}
exec('"""Exact identity audit for the two-boundary entropy lemma.\nProof of strict inequalities and the joint-parameter argument is analytic,\nrecorded separately. This checker is not a numerical certificate.\n"""\nfrom pathlib import Path\nimport json,hashlib\nimport sympy as s\nc,r,q,w,z,x=s.symbols(\'c r q w z x\',real=True)\nd=2*q-1\nchecks=[]\ndef same(name,a,b):\n residual=s.factor(s.together(a-b))\n assert residual==0,(name,residual)\n checks.append(name)\nA1=(d*r+1/r)/2\nsame(\'A1 minus q\',A1-q,(1-r)*(1+(1-2*q)*r)/(2*r))\nsame(\'A_k positivity decomposition at k1\',A1,(1/r-r)/2+q*r)\n# Laurent coefficients yield exact integrals against logarithm of |B|^2.\nB=1-2*c*z+z*z\nBabs=s.expand(B*B.subs(z,1/z))\nS=2+4*c*c\nsame(\'base norm\',Babs.coeff(z,0),S)\nsame(\'base Laurent polynomial\',Babs,S-4*c*(z+1/z)+(z*z+1/z**2))\nL={k:-2*s.chebyshevt(abs(k),c)/abs(k) for k in range(-3,4) if k}\ndef log_pair(A):\n A=s.expand(A)\n return s.expand(sum(A.coeff(z,-k)*L[k] for k in L))\nH0=log_pair(Babs)\nH1=log_pair(s.expand((z+1/z)*Babs/2))\nsame(\'base log integral\',H0,2+12*c*c)\nsame(\'first cosine log integral\',H1,-8*c-s.Rational(8,3)*c**3)\nmean_x=s.expand((z+1/z)*Babs/2).coeff(z,0)/S\nsame(\'mean cosine\',mean_x,-4*c/S)\ncov=s.factor(H1/S-mean_x*H0/S)\nsame(\'covariance formula\',cov,-8*c*(4*c**4-4*c*c+3)/(3*S*S))\nsame(\'positive covariance polynomial\',4*c**4-4*c*c+3,4*(c*c-s.Rational(1,2))**2+2)\nbase_entropy=H0/S-s.log(S)\nexcess=base_entropy-(1-s.log(2))\nsame(\'base entropy excess after w substitution\',s.expand_log(excess.subs(c*c,(w-1)/2),force=True),2-2/w-s.log(w))\nsame(\'base entropy derivative\',s.diff(2-2/w-s.log(w),w),(2-w)/w**2)\nsame(\'log upper comparison derivative\',s.diff((x-1/x)/2-s.log(x),x),(x-1)**2/(2*x*x))\nassert s.Rational(9,16)>s.Rational(1,2)\nchecks.append(\'exact 3/4 > 1/sqrt2 comparison\')\n# Normalized fixed-density logarithmic norm derivative at q=q0.\nq0=s.symbols(\'q0\', real=True)\nJ=s.Function(\'J\')\na=(1-q0)/(1-q)\nN=(1-q)*s.log(J(a))\nres=s.diff(N,q).subs(q,q0).subs(J(1),1)\nsame(\'fixed density parameter derivative at norm1\',res,s.Subs(s.Derivative(J(x),x),x,1))\n', ns)
results.extend('two-boundary' + ': ' + name for name in ns['checks'])

q,z=sp.symbols('q z', real=True)
L=2*z-(1+q)*z**2+(sp.Rational(2,3)+sp.Rational(4,3)*q)*z**3
F=sp.series(sp.exp(L),z,0,4).removeO().expand()
assert sp.simplify(F-(1+2*z+(1-q)*z**2-sp.Rational(2,3)*q*z**3))==0
results.append('canonical first-coefficient function through degree three')
assert sp.binomial(6,3)==20
results.append('additional H1 equality example central coefficient')
report={'status':'PASS_EXACT_IDENTITIES', 'checks_count':len(results),
        'checks':results,
        'scope':'Symbolic identities only. The analytic inequalities, compactness, and global argument are proved in the manuscript.'}
print(json.dumps(report,indent=2))
\end{filecontents*}
\documentclass[11pt]{article}
\usepackage[T1]{fontenc}
\usepackage{lmodern}
\usepackage[a4paper,margin=27mm]{geometry}
\usepackage{amsmath,amssymb,amsthm,mathtools}
\usepackage{booktabs,array,microtype}
\usepackage{xcolor}
\usepackage[colorlinks=true,linkcolor=blue!45!black,citecolor=blue!45!black,urlcolor=blue!45!black]{hyperref}
\usepackage{fancyhdr}
\usepackage{listings}
\usepackage{embedfile}
\pagestyle{fancy}
\fancyhf{}
\fancyhead[L]{\small The third zero-free Hardy coefficient}
\fancyhead[R]{\small Keio}
\fancyfoot[C]{\thepage}
\setlength{\headheight}{14pt}
\setlength{\parskip}{4pt}
\setlength{\parindent}{0pt}
\newtheorem{theorem}{Theorem}[section]
\newtheorem{lemma}[theorem]{Lemma}
\newtheorem{corollary}[theorem]{Corollary}
\newtheorem{proposition}[theorem]{Proposition}
\newcommand{\D}{\mathbb D}
\newcommand{\T}{\mathbb T}
\newcommand{\norm}[1]{\lVert #1\rVert}
\newcommand{\Ent}{\mathcal E}
\newcommand{\Cov}{\operatorname{Cov}}
\newcommand{\Rea}{\operatorname{Re}}
\lstset{basicstyle=\ttfamily\scriptsize,breaklines=true,columns=fullflexible,keepspaces=true,showstringspaces=false,frame=single,rulecolor=\color{black!15}}
\hypersetup{pdftitle={Sharp third coefficients for zero-free Hardy functions},pdfauthor={Keio}}
\begin{document}
\title{Sharp third coefficients for zero-free Hardy functions\\[5pt]
\large An entropy proof of the degree-three\\Hummel--Scheinberg--Zalcman bound}
\author{Keio}
\date{23 November 2026}
\maketitle
\begin{abstract}
We give a proof of the sharp third-coefficient estimate
$|a_3|\le(2/e)^{1-1/p}$ for zero-free functions in the unit ball of $H^p$,
for every $1<p<\infty$, using the established third-coefficient Krzy\.{z}
theorem as the bounded-function input. Multiplicative variations give a
nonnegative trigonometric polynomial $|f|^p$ of degree three. Its spectral
factor and the singular inner measure reduce a maximizer to three
algebraic logarithmic-moment identities. If one or two spectral factors
have boundary zeros, these identities force the normalized boundary
density to have entropy strictly greater than $1-\log2$. If all three do,
there is an exact transfer to the classical bounded problem. Maximizing
the possible violation also over $q=1-1/p$ forces entropy exactly
$1-\log2$, and rules out every stratum. The argument treats the two
parameter endpoints and repeated boundary factors. The finite-$p$
equality cases are the expected rotations of the canonical HSZ function.
No finite numerical search or unrestricted Krzy\.{z} conjecture is used.
\end{abstract}
{\small\setlength{\parskip}{0pt}\tableofcontents}
\section{Statement, normalization, and the established input}
Let $dm=d\theta/(2\pi)$ be normalized arc measure on $\T$ and set
\[
 \|f\|_p=\sup_{0<r<1}\left(\int_{\T}|f(rz)|^p\,dm(z)\right)^{1/p}.
\]
All logarithms are natural. A function is zero-free if it has no zeros
in $\D$; zeros of its boundary values are allowed.

\begin{theorem}\label{thm:main}
Let $1<p<\infty$ and let $f(z)=\sum_{j\ge0}a_jz^j$ be zero-free in $\D$
with $\|f\|_p\le1$. Then
\begin{equation}\label{eq:main}
 |a_3|\le(2/e)^{1-1/p}.
\end{equation}
Equality holds exactly for
\begin{equation}\label{eq:kappa}
 f(z)=\eta\,2^{-1/p}(1+\omega z^3)^{2/p}
 \exp\left[\left(1-\frac1p\right)
             \frac{\omega z^3-1}{\omega z^3+1}\right],
 \qquad |\eta|=|\omega|=1.
\end{equation}
The analytic power is defined by the logarithm of $1+\omega z^3$ that
vanishes at zero. At $p=1$ the sharp coefficient bound is $1$; at
$p=\infty$ it is $2/e$.
\end{theorem}

The established coefficient result used here is the degree-three
Krzy\.{z} theorem: if $F$ is zero-free and $\|F\|_\infty\le1$, then
$|F_3|\le2/e$, with equality precisely for rotations of
$\exp((z^3-1)/(z^3+1))$. The original degree-three result is due to
Hummel, Scheinberg, and Zalcman~\cite{HSZ}; a later independent proof,
including the equality statement, is given by Stupin~\cite{Stupin}.
Only this fixed third-coefficient theorem is invoked, together with
standard Hardy inner--outer factorization, Fej\'er--Riesz factorization,
and elementary compactness and entropy inequalities.

Write
\[
 q=1-1/p\in(0,1),\qquad C=1-\log2,
 \qquad (2/e)^q=e^{-Cq}.
\]
The function in \eqref{eq:kappa} is admissible: its singular exponential
has boundary modulus one almost everywhere, and
\[
 |f|^p=\frac{|1+\omega z^3|^2}{2},\qquad
 \|f\|_p=1,\qquad |a_3|=2^{1-1/p}e^{-q}=e^{-Cq}.
\]
This verifies sharpness and supplies a positive lower bound throughout
the compactness arguments below.

We use two normalizations of the spectral factor: $Q$ has $\|Q\|_2=1$,
whereas $P=Q/Q(0)$ has $P(0)=1$. Let $m$ be the number of factors of $P$
with zeros on $\T$, \emph{counted with multiplicity}. This is distinct
from the number $N$ of positive singular atoms. Zero masses and repeated
boundary factors are included in the arguments.
\section{Existence and first variations}

Existence follows from normal-family compactness and the lower semicontinuity
of the $H^p$ norm: a maximizing sequence has a locally uniform subsequential
limit, and Hurwitz's theorem makes it either zero-free or identically zero.
The explicit HSZ candidate has nonzero third coefficient, so the maximizing
limit is not identically zero. Its norm is one by scaling. Rotation of the
variable and multiplication by a unimodular constant arrange $a_0,a_3>0$.

\begin{lemma}[Boundary polynomial]
Put $w=|f|^p$ on the boundary. Then
\begin{equation}\label{eq:w}
 w(e^{i\theta})=1+2\operatorname{Re}\sum_{k=1}^3
       \frac{a_{3-k}}{a_3}e^{-ik\theta}\geq0.
\end{equation}
In particular, the proposed identity $|f|^p=h/a_3$ is correct with
$h=a_3+2\operatorname{Re}\sum_{k=1}^3 a_{3-k}e^{-ik\theta}$.
\end{lemma}
\begin{proof}
For every analytic polynomial $g$ and real $\epsilon$, the function
$f e^{\epsilon g}/\norm{f e^{\epsilon g}}_p$ is admissible. Differentiating
the real part of its third coefficient at zero gives
$\operatorname{Re}(fg)_3=a_3\int w\operatorname{Re}g\,dm$.
Applying the same assertion to $ig$ yields the complex identity
\begin{equation}\label{eq:stationary}
 (fg)_3=a_3\int_{\T}w g\,dm.
\end{equation}
Taking $g=z^k$ determines every Fourier coefficient of $w$: it is
$a_{3-k}/a_3$ for $1\leq k\leq3$, zero for $k>3$, and its zeroth coefficient
is one. Fourier uniqueness in $L^1$ proves \eqref{eq:w} almost everywhere.
The resulting trigonometric polynomial is nonnegative everywhere by continuity.
\end{proof}

Write the zero-free inner--outer factorization as
\[
 f=\eta I_\mu O,\qquad
 I_\mu(z)=\exp\left[-\int_{\T}\frac{\zeta+z}{\zeta-z}\,d\mu(\zeta)\right],
\]
where $\mu$ is a finite positive singular measure, $O(0)>0$ is outer, and
$|\eta|=1$. Because $f$ has no zeros, there is no Blaschke factor.

\begin{lemma}[Atomic support and exact degree]
The singular measure $\mu$ is supported on the boundary zeros of $w$.
There is an outer polynomial $Q$ of degree exactly three, with $Q(0)>0$
and $\norm Q_2=1$, such that
\[
 w=|Q|^2,\qquad f=\eta I_\mu Q^{2/p}.
\]
Every boundary atom lies at a zero of $Q$; there are at most three distinct
atoms. The converse inclusion of zero sets and atom supports is not asserted.
\end{lemma}
\begin{proof}
Replace $\mu$ by $(1+\epsilon)\mu$ for small real $\epsilon$ of either sign.
This changes only the inner factor, hence leaves the boundary norm one.
Differentiating the third coefficient and taking real parts gives
$\int h\,d\mu=0$. Since $h=a_3w\geq0$, the measure is supported on its zeros.
A nonzero nonnegative trigonometric polynomial of degree three has at most
three distinct zeros, each of even multiplicity. Fej\'er--Riesz factorization
gives $w=|Q|^2$ with $Q$ outer and of degree at most three. The outer factor
of $f$ is $Q^{2/p}$ by uniqueness of outer factors, using the branch positive
at zero. Finally, if $Q=\sum_{j=0}^3q_jz^j$, the highest Fourier coefficient
in \eqref{eq:w} is
$a_0/a_3=q_0\overline{q_3}\ne0$. Consequently $q_3\ne0$ and the degree is three.
\end{proof}

\section{The exact finite stationary equations}
Here and below $\delta=2/p-1=1-2q$.
Write
\[
 Q=S^{-1/2}P,\quad
 P(z)=\prod_{i=1}^3(1-\alpha_i z)=1+u z+v z^2+b z^3,
 \quad S=1+|u|^2+|v|^2+|b|^2.
\]
Every $\alpha_i\ne0$, $|\alpha_i|\leq1$. A boundary root of $Q$ is at
$\overline{\alpha_i}$ when $|\alpha_i|=1$. In this parameter convention write
the positive singular atom masses as $\lambda_j$ at
$\zeta_j=\overline{\alpha_j}$, and set
$t=\sum_j\lambda_j$ and $M_k=\sum_j\lambda_j\alpha_j^k$.
We also write $P(z)=\sum_{i=0}^3p_i z^i$.
The reflected polynomial is $P^*(z)=z^3\overline{P(1/\bar z)}$.
The moment equations below are covariant under a rotation of the variable
and unaffected by a scalar unimodular factor. Thus later arguments may
rotate an atom to $1$ and retain $a_3>0$ without requiring $a_0>0$ in
that same normalization.

\begin{lemma}[Moment equations, including the factor $k$]
The first four Taylor coefficients satisfy
\begin{equation}\label{eq:jet}
 \frac{f(z)}{a_0}=\frac{P(z)P^*(z)}{\overline b}+O(z^4).
\end{equation}
Consequently, the three exact logarithmic equations are
\begin{equation}\label{eq:moments}
 2k M_k=\sum_{i=1}^3\overline{\alpha_i}^{-k}
                  -\delta\sum_{i=1}^3\alpha_i^k,
 \qquad k=1,2,3.
\end{equation}
Moreover,
\begin{equation}\label{eq:value}
 |a_3|=\frac{e^{-t}S^q}{|b|}.
\end{equation}
\end{lemma}
\begin{proof}
The coefficient relation \eqref{eq:stationary} gives
$a_{3-k}/a_3=S^{-1}\sum_{i=0}^{3-k}p_i\bar p_{i+k}$, which is precisely
the coefficient of degree $3-k$ in $PP^*/S$.
The constant coefficient gives $a_0/a_3=\bar b/S$, proving \eqref{eq:jet}.
Also
\[
 \log(f/a_0)=\frac2p\log P-2\sum_{k\geq1}M_k z^k.
\]
Subtracting $\log P$ from the logarithm of \eqref{eq:jet}, and using
\[
 [z^k]\log P=-\frac1k\sum_i\alpha_i^k,
 \qquad
 [z^k]\log(P^*/\bar b)=-\frac1k\sum_i\bar\alpha_i^{-k},
\]
proves \eqref{eq:moments}. Finally $|a_0|=e^{-t}S^{-1/p}$ and
$|a_3/a_0|=S/|b|$, proving \eqref{eq:value}.
\end{proof}

\section{Fixed-inner-factor uniqueness and the absence of the outer stratum}
\begin{lemma}[Fixed-inner-factor maximization]
The stationary function $f$ is the unique normalized maximizer of
$\operatorname{Re}v_3$ over $v\in I_\mu H^p$ with $\norm v_p\leq1$.
Consequently $\mu\ne0$. If $I_\mu$ has real Taylor coefficients, then the
maximizer with $a_3>0$ also has real Taylor coefficients.
\end{lemma}
\begin{proof}
Consider the bounded complex linear functional
\[
 L(v)=v_3-a_3\int_{\T}|f|^{p-2}\bar f v\,dm
\]
on $H^p$. By \eqref{eq:stationary} it vanishes on $f$ times analytic
polynomials. This set is dense in $I_\mu H^p$ by outer cyclicity.
For completeness, in the present polynomial setting no general cyclicity
theorem is needed: with $O=Q^{2/p}$ and $0<r<1$, the functions $1/O(rz)$
are analytic in a neighborhood of the closed disk and are uniform limits of
polynomials. Moreover $O(z)/O(rz)\to1$ almost everywhere and in $L^p$.
Indeed each factor satisfies
$|(1-\alpha z)/(1-r\alpha z)|\leq2$ for $|z|=1$, so dominated convergence
applies to the positive analytic powers of their product. Multiplying by an
arbitrary polynomial shows that $L$ vanishes on $I_\mu$ times polynomials,
and density of polynomials in $H^p$ finishes the argument.
Thus for $v\in I_\mu H^p$, H\"older's inequality gives
$|v_3|\leq a_3\norm v_p$, since $\norm f_p=1$.
Equality in the real functional is unique by the strict convexity of $H^p$.
If $\mu=0$, the subspace is all of $H^p$; the unique coefficient maximizer
is $z^3$, contradicting the absence of zeros in $f$.
If $I_\mu$ is real, $\overline{f(\bar z)}$ lies in the same subspace, has the
same norm, and has the same positive third coefficient. Uniqueness proves
the final assertion.
\end{proof}

\section{Transfer of the full boundary-factor stratum to Krzy\.{z}}
\begin{theorem}[Exact full-boundary transfer]
Suppose all three factors of $P$ have $|\alpha_i|=1$.
After rotations arranging $a_0,a_3>0$, there is a zero-free analytic $F$ with $|F|\leq1$ on $\D$ and
$F_3>0$ such that
\[
 a_3=F_3^q.
\]
In particular, the established third-coefficient Krzy\.{z} inequality
$F_3\leq2/e$ implies the desired HSZ bound $a_3\leq(2/e)^q$
on this entire stratum. Some of its possible singular masses may be zero.
\end{theorem}
\begin{proof}
Normalize $a_0,a_3>0$ as above. Since $|b|=1$ and $a_0/a_3=\bar b/S$,
we have $b=1$. All roots lie on the unit circle, so $P^*=P$.
Equation \eqref{eq:jet} becomes $f/a_0=P^2+O(z^4)$.
Let $F=I_{\mu/q}$. This is a zero-free function bounded by one, because
$q>0$ and $\mu$ is positive. Comparing logarithms gives
$F/F_0=P^2+O(z^4)$.
As $P^*=P$, one has $[z^3]P^2=\sum_{i=0}^3|p_i|^2=S$.
Thus $F_3=e^{-t/q}S$, whereas \eqref{eq:value} is $a_3=e^{-t}S^q$.
The claimed identity follows. No limiting argument at $p=1$ is used.
\end{proof}


\section{Entropy under a positive multiplicative tilt}
For a probability density $W$ relative to $dm$, put
\[
 \Ent(W)=\int_{\T}W\log W\,dm,\qquad 0\log0=0.
\]
If $W_0dm=\nu$ is a probability measure and $J$ is positive and bounded,
let $W=W_0J/\mathbb E_\nu J$. Direct expansion gives
\begin{equation}\label{eq:tilt}
 \Ent(W)=D(Wdm\|\nu)+\Ent(W_0)
 +\frac{\Cov_\nu(J,\log W_0)}{\mathbb E_\nu J},
\end{equation}
where $D(\rho\|\nu)=\int\log(d\rho/d\nu)\,d\rho\ge0$ is relative entropy.
It is strictly positive unless $J$ is constant $\nu$-almost everywhere.
In our applications $J$ is bounded away from zero and the zeros of $W_0$
are finite, so every term is integrable. Also, for two real functions $A,B$,
\[
 \Cov_\nu(A,B)=\frac12\iint(A(x)-A(y))(B(x)-B(y))\,d\nu(x)d\nu(y).
\]
In particular it is positive for two strictly monotone functions with
the same direction on a nondegenerate distribution.
\section{One boundary factor: strict entropy}
\begin{lemma}\label{lem:one}
At a third-coefficient extremizer with exactly one boundary factor of
$P$, the density $W=|P|^2/\|P\|_2^2$ satisfies $\Ent(W)>C$.
\end{lemma}
\begin{proof}
Let $0<q<1$, $d=2q-1$, and suppose a stationary candidate has one boundary factor, with a positive singular atom there. Rotate the sole singular atom to $1$ and choose $a_3>0$.
Fixed-inner-factor uniqueness makes $f$ real; the unique outer spectral
factor normalized at zero is then real as well. Thus write
\[
 P(z)=(1-z)(1-sz+vz^2),\qquad s,v\in\mathbb R,
\]
where the reciprocal roots of $1-sz+vz^2$ lie strictly in $\mathbb D$. Since $P$ has degree three, $v\ne0$. For a real quadratic with reciprocal roots strictly in $\D$, the
conditions are
\[
 -1<v<1,\qquad |s|<1+v.
\]
Let $\lambda>0$ denote the singular mass. The first two log-jet identities are
\[
 2\lambda=2q+s(d+v^{-1}),\qquad
 4\lambda=2q+(s^2-2v)(d+v^{-2}).
\]
Consequently
\[
 E(q;s,v):=(s^2-2v)(d v^2+1)-2s(dv+1)v-2qv^2=0.
\]
This expression is affine in $q$.

\paragraph{Step 1: $v>0$.}
If $v=-u$ with $0<u<1$, let $a=1-u$; stability gives $-a<s<a$. At $q=0$,
\[
 E(0;s,-u)=(1+u)\bigl(a s^2+2u(s+a)\bigr)>0.
\]
At $q=1$,
\[
 E(1;s,-u)=(1+u^2)s^2+2ua s+2u(1-u+u^2)>2u^2>0,
\]
Indeed the difference from $2u^2$ is
$(1+u^2)s^2+2u(1-u)(s+1-u)>0$. Affineness contradicts $E(q)=0$.

\paragraph{Step 2: $s>0$.}
Assume $v>0$ and $s=-x\le0$. The first moment and $\lambda>0$ give
\[
 0<\lambda=q(1-x)-\frac{x(1-v)}{2v}.
\]
Thus $x<1$ and, since $q<1$,
\[
 0\le x<\frac{2v}{1+v}=:X.
\]
Both endpoint expressions below are increasing functions of $x\ge0$:
\[
 E(0;-x,v)=(1-v)\bigl((1+v)x^2+2vx-2v(1+v)\bigr),
\]
\[
 E(1;-x,v)=(1+v^2)x^2+2v(1+v)x-2v(1+v+v^2).
\]
At $x=X$ they equal, respectively,
\[
 -\frac{2v(1-v)^3}{1+v}<0,
 \qquad
 -\frac{2v(1-v)^2(v^2+v+1)}{(1+v)^2}<0.
\]
Therefore $E(q)<0$, again a contradiction.

\paragraph{Step 3: a decreasing outer tilt.}
For $s>0$ and $0<v<1$, stability gives $s<1+v<2$. Put
\[
 A=(1+v)s^2-2vs-2v(1+v),\qquad
 B=s^2-2s-2v-1<0.
\]
Then $E(q)=(1-v)A+2qv^2B=0$, so $A>0$. The quadratic $A$ has one positive root, and at $s=4v/(1+v)$ its value is
\[
 -\frac{2v(1-v)^2}{1+v}<0.
\]
It follows that
\[
 s>\frac{4v}{1+v}.
\]
For $x=\cos\theta$,
\[
 J(x):=|1-se^{i\theta}+ve^{2i\theta}|^2
 =(1-v)^2+s^2-2s(1+v)x+4vx^2.
\]
Hence $J'(x)=-2s(1+v)+8vx<0$ on $[-1,1]$.

\paragraph{Step 4: entropy.}
Let $dm=d\theta/(2\pi)$, $W_0=|1-e^{i\theta}|^2/2=1-\cos\theta$, and let $d\nu=W_0dm$. With $S=\|P\|_2^2$,
\[
 W=|P|^2/S=W_0J/\mathbb E_\nu J.
\]
Relative entropy and the covariance inequality give
\begin{align*}
 \int W\log W\,dm
 &=D(Wdm\|\nu)+\mathbb E_{Wdm}\log W_0\\
 &\ge \mathbb E_\nu\log W_0+
 \frac{\operatorname{Cov}_\nu(J,\log W_0)}{\mathbb E_\nu J}
 >1-\log2.
\end{align*}
The covariance is strictly positive because both $J(x)$ and $\log(1-x)$ are strictly decreasing and $\nu$ is not a point mass. Strictness also follows directly from the covariance double-integral
identity above. The endpoint integral is
$\int (1-\cos\theta)\log(1-\cos\theta)\,dm=1-\log2$.

 Its value follows from the constant and first cosine
coefficients of the Abel-summed $L^1$ Fourier identity
$\log(1-\cos\theta)=-\log2-2\sum_{k\ge1}\cos(k\theta)/k$.
\end{proof}
\section{Two boundary factors: alignment and strict entropy}
\begin{lemma}\label{lem:two}
At a third-coefficient extremizer with exactly two boundary factors of
$P$, counted with multiplicity, the density $W=|P|^2/\|P\|_2^2$
satisfies $\Ent(W)>C$.
\end{lemma}
\begin{proof}
Let $0<q<1$, let $d=2q-1$, and suppose a stationary cubic spectral polynomial has two boundary reciprocal roots (counted with multiplicity) and one reciprocal root strictly inside the disk. Rotate so that
\[
 \begin{gathered}
 P(z)=B(z)(1-\beta z),\qquad B(z)=1-2cz+z^2,\\
 c=\cos\theta\in[-1,1],\qquad\beta=re^{i\phi},\quad 0<r<1.
 \end{gathered}
\]
The two boundary reciprocal roots are $e^{\pm i\theta}$. Let their nonnegative singular masses be $\lambda_1,\lambda_2$, allowing a zero mass. At a repeated boundary root only their sum matters. Write $t=\lambda_1+\lambda_2$ and $\Delta=\lambda_1-\lambda_2$.
The correct three log-jet equations are
\[
 (kt-2q)T_k(c)+ik\Delta\sin(k\theta)=A_k e^{ik\phi},\qquad k=1,2,3,
\]
where
\[
 A_k=\frac{d r^k+r^{-k}}2>0.
\]
In addition,
\[
 A_1-q=\frac{(1-r)(1+(1-2q)r)}{2r}>0.
\]

\paragraph{Step 1: alignment, $c\operatorname{Re}\beta\ge0$.}
If $\sin\phi\ne0$, the imaginary parts for $k=1,2$ give
\[
 \Delta\sin\theta=A_1\sin\phi,
 \qquad 4\Delta c\sin\theta=2A_2\sin\phi\cos\phi.
\]
Consequently
\[
 \cos\phi=\frac{2cA_1}{A_2},\qquad
 c\operatorname{Re}\beta=\frac{2r c^2 A_1}{A_2}\ge0.
\]
This branch automatically has $\sin\theta\ne0$.

If $\sin\phi=0$, change $z$ to $-z$ if necessary to make $\beta=r>0$; the product $c\operatorname{Re}\beta$ is unchanged by this operation. The first real equation is $(t-2q)c=A_1$, so $c\ne0$. Suppose $c=-u<0$. Then
\[
 t=2q-A_1/u,\qquad
 2t-2q=2(q-A_1/u)<0,
\]
because $A_1>q$ and $0<u\le1$. The second real equation forces
\[
 2u^2-1<0,\qquad u<1/\sqrt2.
\]
Thus $T_3(-u)=u(3-4u^2)>0$. The third real equation forces $3t-2q>0$, hence
\[
 4q-3A_1/u>0,\qquad
 u>\frac{3A_1}{4q}>\frac34>\frac1{\sqrt2},
\]
a contradiction. Thus $c>0$ and alignment holds. This reasoning also covers repeated boundary roots $c=\pm1$; in that case the first imaginary equation already forces $\sin\phi=0$.

\paragraph{Step 2: the quadratic boundary density.}
Let $dm=du/(2\pi)$, put $x=\cos u$, and write
\[
 S_B=\|B\|_2^2=2+4c^2,\qquad W_B=|B|^2/S_B,\qquad d\nu=W_Bdm.
\]
The Fourier expansion of the logarithms of the two unit-circle factors, paired with the finite Fourier polynomials $|B|^2$ and $x|B|^2$, gives
the following identities. These pairings can be justified by radial Abel
limits in $L^1$; only the first three Fourier modes occur.
\[
 \int |B|^2\log|B|^2\,dm=2+12c^2,
 \qquad
 \int x|B|^2\log|B|^2\,dm=-8c-\frac83c^3.
\]
Also $\int x|B|^2dm=-4c$. Therefore
\[
 \operatorname{Cov}_{\nu}(x,\log W_B)
 =-\frac{8c(4c^4-4c^2+3)}{3S_B^2}.
\]
The factor $4c^4-4c^2+3=4(c^2-1/2)^2+2$ is positive.
The entropy of the base density is
\[
 \mathcal E(W_B)=\frac{1+6c^2}{1+2c^2}-\log(2+4c^2)\ge1-\log2.
\]
Indeed, with $w=1+2c^2\in[1,3]$, its excess over $1-\log2$ is
$2-2/w-\log w$. This function increases up to $w=2$ and then decreases; its endpoint values are $0$ and $4/3-\log3>0$.
For the last strict inequality, differentiate
$(x-x^{-1})/2-\log x$: its derivative is $(x-1)^2/(2x^2)>0$
for $x>1$, and evaluate at $x=3$. Thus equality holds only for $c=0$.

\paragraph{Step 3: the inside-root tilt raises entropy.}
Let
\[
 J(u)=|1-\beta e^{iu}|^2
 =1+r^2-2r\cos\phi\cos u+2r\sin\phi\sin u.
\]
The density $W_B$ and its logarithm are even in $u$, so the odd term contributes zero covariance. By Step 1 and the covariance identity,
\[
 \operatorname{Cov}_{\nu}(J,\log W_B)
 =-2r\cos\phi\operatorname{Cov}_{\nu}(x,\log W_B)\ge0.
\]
For $W=|P|^2/\|P\|_2^2=W_BJ/\mathbb E_\nu J$,
\begin{align*}
 \mathcal E(W)
 &=D(Wdm\|\nu)+\mathcal E(W_B)
   +\frac{\operatorname{Cov}_\nu(J,\log W_B)}{\mathbb E_\nu J}\\
 &>\mathcal E(W_B)\ge1-\log2.
\end{align*}
Strictness holds because $r>0$ makes $J$ nonconstant on the full-measure support of $\nu$, so $D(Wdm\|\nu)>0$.

\end{proof}
\section{Maximizing also over the Hardy exponent}\label{sec:joint}
Put $q=1-1/p$, $C=1-\log2$, and let $M(q)$ be the supremum of the
normalized third coefficient over zero-free $H^{1/(1-q)}$ functions.
For $0<q<1$ the stationary reduction proves attainment and supplies
$f_q=\eta I_\mu Q^{2/p}$ with $\|Q\|_2=1$, $\deg Q=3$.
Consequently $\|Q\|_\infty\leq2$ and
\[
 \|f_q\|_\infty\leq4^{1/p}=4^{1-q}.
\]
Applying the established third-coefficient Krzy\.{z} inequality to
$f_q/\|f_q\|_\infty$ gives
\[
 (2/e)^q\leq M(q)\leq (2/e)4^{1-q}.
\]
Here the lower bound is provided by the explicit HSZ candidate.
Thus $M(q)\to2/e$ as $q\uparrow1$. At the other endpoint,
$(2/e)^q\leq M(q)\leq1$, so $M(q)\to1$ as $q\downarrow0$.

The function $M$ is continuous on $(0,1)$. One direct proof compares a
normalized extremizer at either of two nearby parameter values. If its
own exponent is $p_0$, its boundary density is
$W=|f|^{p_0}=|Q|^2$, with $0\leq W\leq4$ and $\int W=1$.
At another exponent $p$, the logarithmic norm is
\[
 \log\|f\|_p=\frac1p\log\int W^{p/p_0}\,dm.
\]
For $p,p_0$ in a fixed compact subset of $(1,\infty)$ and sufficiently
close, the first two parameter derivatives of this expression are bounded
uniformly in $W$: the functions $x^a|\log x|^j$ are bounded on $[0,4]$
for $a\geq1/2$, $j\leq2$, and the integrals $\int W^a$ are bounded below
uniformly by $\min(1,4^{a-1})$.
Using the extremizer at either parameter gives local Lipschitz bounds for
$\log M$. The preceding endpoint limits extend $M$ continuously to $[0,1]$.

Suppose the HSZ bound fails somewhere. Then
\[
 R(q)=e^{Cq}M(q)
\]
has a maximum greater than one at an interior $q_*\in(0,1)$.
Let $f_*$ be a normalized extremizer at $p_*=1/(1-q_*)$.
This function is bounded and zero-free, so for every $q\in(0,1)$ it is an
admissible test function after normalizing its $H^{1/(1-q)}$ norm. Hence
\[
 G(q)=e^{Cq}\frac{|(f_*)_3|}{\|f_*\|_{1/(1-q)}}\leq R(q)\leq R(q_*),
 \qquad G(q_*)=R(q_*).
\]
The ordinary derivative $G'(q_*)$ vanishes. Differentiation of the fixed
function's norm gives
\[
 \left.\frac{d}{dq}\log\|f_*\|_{1/(1-q)}\right|_{q=q_*}
 =\int_{\T}W\log W\,dm,
 \qquad W=|f_*|^{p_*},\quad\int W=1.
\]
Therefore every joint-parameter maximizing counterexample obeys the exact
additional equation
\[
 \boxed{\int W\log W\,dm=1-\log2.}
\]
This is a condition at a maximizer over the parameter $q$, not a condition
that may be imposed on every fixed-$p$ extremizer. No differentiability of
the maximizing function $f_q$ or of the envelope $M(q)$ is needed.


\section{Proof of the global inequality and the equality cases}
\begin{proof}[Proof of Theorem~\ref{thm:main}]
Suppose first that \eqref{eq:main} fails. Section~\ref{sec:joint} gives an
interior joint maximizing counterexample $f_*$ with
$\Ent(|f_*|^{p_*})=C$. Its normalized spectral polynomial has degree three.
There is at least one boundary factor because its singular measure is
nonzero and supported on boundary zeros. Count boundary factors with
multiplicity. If $m=1$, Lemma~\ref{lem:one} gives entropy strictly greater
than $C$. If $m=2$, Lemma~\ref{lem:two} gives the same strict inequality,
including a possible repeated factor. If $m=3$, the exact full-boundary
transfer already gives $|a_3|\le(2/e)^{q_*}$.
Each case contradicts the defining properties of the joint maximum.
This proves the inequality for every $1<p<\infty$.

Now let $f$ attain equality at a fixed $q_0\in(0,1)$, with norm one.
It is a bounded stationary extremizer. Since the inequality has now been
proved for \emph{all} $q$, the fixed-function ratio
\[
 G(q)=e^{Cq}\frac{|a_3(f)|}{\|f\|_{1/(1-q)}}
\]
is at most one everywhere and equals one at $q_0$. Its ordinary derivative
therefore gives $\Ent(|f|^{p_0})=C$. The two strict entropy lemmas exclude
$m=1,2$, so all three boundary factors are present.

The transfer $F=I_{\mu/q_0}$ satisfies
$|F_3|^{q_0}=|a_3(f)|=(2/e)^{q_0}$. The classical degree-three Krzy\.{z}
equality classification~\cite{Stupin,HSZ}, together with $F(0)>0$, yields
\[
 F(z)=\exp\left(\frac{\omega z^3-1}{\omega z^3+1}\right),\qquad|\omega|=1.
\]
Its singular measure consists of three equal positive atoms. Their support
forces $P(z)=1+\omega z^3$ and $S=2$, and $I_\mu=F^{q_0}$.
Substitution in $f=\eta I_\mu P^{2/p_0}/S^{1/p_0}$ proves
\eqref{eq:kappa}. The direct verification in Section~1 proves the converse.
The endpoint inequalities are the elementary $H^1$ coefficient estimate
and the established third-coefficient Krzy\.{z} theorem, respectively.
\end{proof}

\section{Endpoints, scope, and exact verification}
The parameter argument does not assume a positive margin uniform as
$p\downarrow1$. Instead the exact endpoint limits put a hypothetical
strict violation at an interior parameter, where the entropy derivative
is legitimate. Equality at $p=1$ is not uniquely canonical: both
\[
 \frac{(1+z^3)^2}{2}\quad\hbox{and}\quad\frac{(1+z)^6}{20}
\]
are zero-free in $\D$, have $H^1$ norm one and third coefficient one.
For the second example, $\int|1+z|^6dm=\binom63=20$.

The proof is a degree-three argument. Its one-factor algebra and
two-factor alignment use the first three logarithmic moments in an
essential way. No all-degree Hardy coefficient theorem or all-degree
Krzy\.{z} conjecture is inferred.

The accompanying exact symbolic script checks the displayed one-factor
polynomial identities, two-factor Fourier pairings and covariance,
elementary derivative identities, and a canonical Taylor expansion.
It uses SymPy and performs no numerical search. Strict positivity,
compactness, relative entropy, and the global contradiction are proved
analytically above; passing the script is not a formal proof of those
steps. The TeX source contains the whole script and writes it on
compilation. It is also attached to the PDF.
\begin{lstlisting}
python verify_hsz_third_coefficient.py
\end{lstlisting}
\embedfile[desc={Exact symbolic identity checks for the third-coefficient proof},ucfilespec={verify\string_hsz\string_third\string_coefficient.py}]{verify_hsz_third_coefficient.py}

\appendix
\section{A distinction from a claimed ambient-openness argument}
The proof above uses multiplicative variations that preserve absence of
zeros. It does not need the set of all zero-free Hardy functions to be
open in the surrounding Hardy norm. In fact that ambient-openness
statement, appearing as Lemma~3 in~\cite{Krushkal2022}, is false as written.
Let
\[
 f(z)=\frac{1-z}{4},\qquad
 g_\varepsilon(z)=\frac{1-z}{4}-\varepsilon,
 \qquad 0<\varepsilon<\frac14.
\]
The function $f$ is zero-free in $\D$, $\|f\|_p\le1/2$, and
$\|g_\varepsilon-f\|_p=\varepsilon\to0$. Nevertheless
$g_\varepsilon(1-4\varepsilon)=0$, while
$\|g_\varepsilon\|_p\le1/2+\varepsilon<1$.
This observation concerns that stated lemma, not a counterexample to
the HSZ inequality. The later general comparison in~\cite{Krushkal2026}
is not an input to the present proof.

\begin{thebibliography}{9}
\bibitem{HSZ} J. A. Hummel, S. Scheinberg, and L. Zalcman,
\emph{A coefficient problem for bounded nonvanishing functions},
Journal d'Analyse Math\'ematique \textbf{31} (1977), 169--190.
\href{https://doi.org/10.1007/BF02813302}{doi:10.1007/BF02813302}.
\bibitem{Stupin} D. L. Stupin,
\emph{A new proof of the Krzyz conjecture for $n=3$},
Izvestiya of Saratov University. Mathematics. Mechanics. Informatics
\textbf{24}(3) (2024), 342--350, especially Corollary~1 on p.~349.
\href{https://doi.org/10.18500/1816-9791-2024-24-3-342-350}{doi:10.18500/1816-9791-2024-24-3-342-350}.
\bibitem{Krushkal2022} S. L. Krushkal,
\emph{Quasiconformal deformations of nonvanishing $H^p$ functions
and the Hummel--Scheinberg--Zalcman conjecture},
\href{https://arxiv.org/abs/2211.16295}{arXiv:2211.16295v1} (2022).
\bibitem{Krushkal2026} S. L. Krushkal,
\emph{A New Method for Estimating the Coefficients of Holomorphic Functions},
Axioms \textbf{15}(5) (2026), 361.
\href{https://doi.org/10.3390/axioms15050361}{doi:10.3390/axioms15050361}.
\end{thebibliography}
\end{document}
